\documentclass[12pt, a4paper]{ctexart}
\usepackage{amsmath, amssymb, amsthm}
\usepackage{geometry}
\usepackage{fancyhdr}
\usepackage{enumitem}
\usepackage{tikz}
\usepackage{pgfplots}
\usepackage{array}
\usepackage{multirow}
\usepackage{booktabs}

\geometry{left=2cm, right=2cm, top=2.5cm, bottom=2.5cm}

\setCJKmainfont{Hiragino Sans GB}
\setCJKsansfont{Heiti SC}

\pagestyle{fancy}
\fancyhf{}
\fancyhead[L]{空间向量与立体几何 测试卷}
\fancyhead[R]{数学（选择性必修第一册）}
\fancyfoot[C]{第 \thepage\ 页}
\renewcommand{\headrulewidth}{0.4pt}

\setlength{\parindent}{0pt}

\newcommand{\score}[1]{\hfill\textbf{（#1分）}}

\begin{document}

\begin{center}
{\Large\bfseries 空间向量与立体几何 测试卷}\\[0.5em]
{\normalsize 时间：120分钟 \quad 满分：150分}
\end{center}

\vspace{1em}

\noindent\textbf{一、单项选择题（本大题共8小题，每小题5分，共40分。每小题只有一个选项符合题目要求。）}

\vspace{0.5em}

\begin{enumerate}
\item 已知空间向量 $\vec{a}=(1,2,3)$，$\vec{b}=(2,-1,1)$，则 $\vec{a}\cdot\vec{b}=$ \score{5}
\begin{enumerate}[label=\Alph*.]
\item $3$
\item $4$
\item $5$
\item $6$
\end{enumerate}

\item 已知向量 $\vec{a}=(1,0,-1)$，$\vec{b}=(0,1,1)$，则 $\vec{a}\times\vec{b}=$ \score{5}
\begin{enumerate}[label=\Alph*.]
\item $(1,-1,1)$
\item $(1,1,-1)$
\item $(-1,1,1)$
\item $(1,-1,-1)$
\end{enumerate}

\item 在空间直角坐标系中，点 $A(1,2,3)$，$B(4,0,5)$，则 $|\overrightarrow{AB}|=$ \score{5}
\begin{enumerate}[label=\Alph*.]
\item $\sqrt{14}$
\item $\sqrt{17}$
\item $\sqrt{21}$
\item $\sqrt{23}$
\end{enumerate}

\item 已知空间向量 $\vec{a}=(2,-1,2)$，$\vec{b}=(4,2,x)$，若 $\vec{a}\perp\vec{b}$，则 $x$ 的值为 \score{5}
\begin{enumerate}[label=\Alph*.]
\item $-3$
\item $3$
\item $-2$
\item $2$
\end{enumerate}

\item 已知平面 $\alpha$ 的一个法向量为 $\vec{n}=(1,-2,2)$，直线 $l$ 的方向向量为 $\vec{v}=(2,3,-1)$，则直线 $l$ 与平面 $\alpha$ 所成角的正弦值为 \score{5}
\begin{enumerate}[label=\Alph*.]
\item $\dfrac{8}{3\sqrt{39}}$
\item $\dfrac{8}{3\sqrt{13}}$
\item $\dfrac{4}{3\sqrt{13}}$
\item $\dfrac{2}{3\sqrt{13}}$
\end{enumerate}

\item 在正方体 $ABCD\text{-}A_1B_1C_1D_1$ 中，棱长为 $2$，以 $D$ 为原点，$\overrightarrow{DA}$，$\overrightarrow{DC}$，$\overrightarrow{DD_1}$ 的方向分别为 $x$，$y$，$z$ 轴的正方向建立空间直角坐标系，则异面直线 $A_1B$ 与 $B_1D_1$ 所成角的余弦值为 \score{5}
\begin{enumerate}[label=\Alph*.]
\item $\dfrac{1}{3}$
\item $\dfrac{1}{2}$
\item $\dfrac{\sqrt{2}}{2}$
\item $\dfrac{\sqrt{3}}{3}$
\end{enumerate}

\item 已知三点 $A(1,0,0)$，$B(0,1,0)$，$C(0,0,1)$，则平面 $ABC$ 的方程为 \score{5}
\begin{enumerate}[label=\Alph*.]
\item $x+y+z=0$
\item $x+y+z=1$
\item $x-y+z=1$
\item $x+y-z=1$
\end{enumerate}

\item 已知 $\vec{a}=(1,1,0)$，$\vec{b}=(1,0,1)$，$\vec{c}=(0,1,1)$，则以 $\vec{a}$，$\vec{b}$，$\vec{c}$ 为棱的平行六面体的体积为 \score{5}
\begin{enumerate}[label=\Alph*.]
\item $1$
\item $2$
\item $3$
\item $4$
\end{enumerate}
\end{enumerate}

\vspace{1em}

\noindent\textbf{二、多项选择题（本大题共4小题，每小题5分，共20分。每小题有多个选项符合题目要求，全部选对得5分，部分选对得2分，有选错的得0分。）}

\vspace{0.5em}

\begin{enumerate}[resume]
\item 已知空间向量 $\vec{a}=(1,2,-2)$，$\vec{b}=(-2,-4,4)$，则下列结论正确的是 \score{5}
\begin{enumerate}[label=\Alph*.]
\item $\vec{a}$ 与 $\vec{b}$ 共线
\item $|\vec{a}|=3$
\item $\vec{a}\cdot\vec{b}=-18$
\item $\vec{a}$ 与 $\vec{b}$ 的夹角为 $0°$
\end{enumerate}

\item 在空间直角坐标系中，已知点 $P(1,1,1)$，$A(2,0,0)$，$B(0,2,0)$，$C(0,0,2)$，则下列结论正确的是 \score{5}
\begin{enumerate}[label=\Alph*.]
\item $\overrightarrow{PA}\perp\overrightarrow{PB}$
\item $\overrightarrow{PA}\perp\overrightarrow{PC}$
\item $\overrightarrow{PB}\perp\overrightarrow{PC}$
\item 点 $P$ 是 $\triangle ABC$ 的外心
\end{enumerate}

\item 已知平面 $\alpha$ 经过点 $A(1,0,0)$，$B(0,1,0)$，$C(0,0,1)$，则下列向量中是平面 $\alpha$ 的法向量的有 \score{5}
\begin{enumerate}[label=\Alph*.]
\item $(1,1,1)$
\item $(-1,-1,-1)$
\item $(2,2,2)$
\item $(1,-1,1)$
\end{enumerate}

\item 关于空间向量，下列说法正确的是 \score{5}
\begin{enumerate}[label=\Alph*.]
\item 若 $\vec{a}\cdot\vec{b}=0$，则 $\vec{a}\perp\vec{b}$
\item 若 $\vec{a}\times\vec{b}=\vec{0}$，则 $\vec{a}$ 与 $\vec{b}$ 共线
\item $|\vec{a}\times\vec{b}|=|\vec{a}||\vec{b}|\sin\theta$，其中 $\theta$ 为 $\vec{a}$ 与 $\vec{b}$ 的夹角
\item 若 $\vec{a}$，$\vec{b}$，$\vec{c}$ 共面，则 $\vec{a}\cdot(\vec{b}\times\vec{c})=0$
\end{enumerate}
\end{enumerate}

\vspace{1em}

\noindent\textbf{三、填空题（本大题共4小题，每小题5分，共20分。）}

\vspace{0.5em}

\begin{enumerate}[resume]
\item 已知空间向量 $\vec{a}=(1,2,2)$，$\vec{b}=(2,-2,1)$，则 $\cos\langle\vec{a},\vec{b}\rangle=$ \underline{\hspace{3cm}}。\score{5}

\item 已知点 $A(1,2,3)$，$B(3,4,5)$，点 $P$ 在线段 $AB$ 上，且 $\overrightarrow{AP}=2\overrightarrow{PB}$，则点 $P$ 的坐标为 \underline{\hspace{3cm}}。\score{5}

\item 在正四面体 $ABCD$ 中，棱长为 $2$，以 $A$ 为原点建立空间直角坐标系，使 $B$ 在 $x$ 轴正半轴上，$C$ 在 $xOy$ 平面第一象限内，则 $D$ 的坐标为 \underline{\hspace{3cm}}。\score{5}

\item 已知平面 $\alpha$ 过点 $(1,0,0)$，且法向量为 $\vec{n}=(2,-1,3)$，点 $P(3,1,2)$ 到平面 $\alpha$ 的距离为 \underline{\hspace{3cm}}。\score{5}
\end{enumerate}

\vspace{1em}

\noindent\textbf{四、解答题（本大题共6小题，共70分。解答应写出文字说明、证明过程或演算步骤。）}

\vspace{0.5em}

\begin{enumerate}[resume]
\item （10分）已知空间向量 $\vec{a}=(1,1,0)$，$\vec{b}=(1,0,1)$，$\vec{c}=(0,1,1)$。
\begin{enumerate}[label=(\arabic*)]
\item 求 $\vec{a}\cdot\vec{b}$，$\vec{b}\cdot\vec{c}$，$\vec{c}\cdot\vec{a}$；
\item 求 $|2\vec{a}-\vec{b}+\vec{c}|$。
\end{enumerate}

\item （10分）在正方体 $ABCD\text{-}A_1B_1C_1D_1$ 中，棱长为 $1$，以 $D$ 为原点，$\overrightarrow{DA}$，$\overrightarrow{DC}$，$\overrightarrow{DD_1}$ 的方向分别为 $x$，$y$，$z$ 轴的正方向建立空间直角坐标系。
\begin{enumerate}[label=(\arabic*)]
\item 写出各顶点的坐标；
\item 求 $A_1C$ 的中点坐标；
\item 求二面角 $A\text{-}BD_1\text{-}C$ 的余弦值。
\end{enumerate}

\item （12分）已知四面体 $OABC$ 中，$\overrightarrow{OA}=\vec{a}$，$\overrightarrow{OB}=\vec{b}$，$\overrightarrow{OC}=\vec{c}$，$M$，$N$ 分别是 $OA$，$BC$ 的中点。
\begin{enumerate}[label=(\arabic*)]
\item 用 $\vec{a}$，$\vec{b}$，$\vec{c}$ 表示 $\overrightarrow{MN}$；
\item 若 $|\vec{a}|=|\vec{b}|=|\vec{c}|=1$，$\vec{a}\cdot\vec{b}=\vec{b}\cdot\vec{c}=\vec{c}\cdot\vec{a}=0$，求 $|\overrightarrow{MN}|$。
\end{enumerate}

\item （12分）如图，在三棱柱 $ABC\text{-}A_1B_1C_1$ 中，侧面 $ABB_1A_1\perp$ 底面 $ABC$，$AB=BC=2$，$\angle ABC=90°$，$AA_1=2$，$E$ 为 $BB_1$ 的中点。
\begin{enumerate}[label=(\arabic*)]
\item 证明：$A_1E\perp AC$；
\item 求直线 $A_1E$ 与平面 $ACC_1A_1$ 所成角的正弦值。
\end{enumerate}

\item （12分）在三棱锥 $P\text{-}ABC$ 中，$PA\perp$ 平面 $ABC$，$AB\perp BC$，$PA=AB=BC=2$。
\begin{enumerate}[label=(\arabic*)]
\item 建立适当的空间直角坐标系，写出各顶点的坐标；
\item 求二面角 $P\text{-}BC\text{-}A$ 的大小；
\item 求点 $A$ 到平面 $PBC$ 的距离。
\end{enumerate}

\item （14分）如图，在四棱锥 $P\text{-}ABCD$ 中，底面 $ABCD$ 是正方形，$PA\perp$ 底面 $ABCD$，$PA=AB=2$，$E$，$F$ 分别是 $PB$，$PD$ 的中点。
\begin{enumerate}[label=(\arabic*)]
\item 证明：$EF\parallel$ 平面 $ABCD$；
\item 证明：$BD\perp EF$；
\item 求二面角 $P\text{-}BD\text{-}A$ 的正切值；
\item 求三棱锥 $P\text{-}AEF$ 的体积。
\end{enumerate}
\end{enumerate}

\end{document}
